% =====================================================================
%  Packages
% =====================================================================

% ----- Langue & encodage ---------------------------------------------
\usepackage[utf8]{inputenc}
\usepackage[T1]{fontenc}
\usepackage[french]{babel}
\usepackage{csquotes}                 % requis par biblatex avec babel

% ----- Police Times-like (substitut libre de Times New Roman) --------
\usepackage{newtxtext}                % corps de texte
\usepackage{newtxmath}                % équations cohérentes avec Times

% ----- Mise en page --------------------------------------------------
\usepackage[a4paper,margin=25mm]{geometry}   % marges 25 mm (Sayadi)
\usepackage{setspace}                         % interligne 1.5
\usepackage{indentfirst}                      % premier paragraphe indenté
\usepackage{parskip}                          % petite espace entre paragraphes

% ----- Graphiques ----------------------------------------------------
\usepackage{graphicx}
\graphicspath{{figures/}}
\usepackage{tikz}                            % placeholders figures
\usetikzlibrary{positioning,arrows.meta,shapes.geometric,calc}
\usepackage{float}                            % option [H] pour figures/tableaux

% ----- Tableaux ------------------------------------------------------
\usepackage{booktabs}                         % règles propres : \toprule \midrule
\usepackage{array}                            % colonnes avancées
\usepackage{tabularx}                         % colonnes auto-ajustées
\usepackage{multirow}
\usepackage{caption}
\captionsetup[figure]{labelfont=bf,font=small}
\captionsetup[table]{labelfont=bf,font=small,position=above}

% ----- Mathématiques -------------------------------------------------
\usepackage{amsmath,amsfonts}
\let\Bbbk\relax                               % évite le conflit avec newtxmath
\usepackage{amssymb}
\usepackage{mathtools}
\usepackage{siunitx}                          % unités cohérentes (\SI{0.95}{})

% ----- Algorithmes (encadrés avec titre, sans \fbox) ------------------
\usepackage{algorithm}
\usepackage{algpseudocode}

% ----- Hyperliens ----------------------------------------------------
\usepackage[hidelinks,breaklinks=true]{hyperref}
\usepackage{url}

% ----- Bibliographie (style IEEE, ordre de citation) -----------------
\usepackage[
  backend=biber,
  style=ieee,
  sorting=none,
  doi=false,isbn=false,url=false
]{biblatex}

% ----- En-tête / pied de page ----------------------------------------
\usepackage{fancyhdr}

% ----- Listes & divers -----------------------------------------------
\usepackage{enumitem}
\usepackage{xcolor}
\usepackage{lipsum}                           % texte fictif — à retirer en production
